% Clase 17 --- Por qué en el centro de la espira el campo NO se anula (§1.4).
% "Recuerden el caso eléctrico... teníamos un aro cargado. ¿Cómo era el campo
%  eléctrico en el centro? Cero. Fíjense que acá no pasa eso."
% Dos paneles con la misma geometría (el aro visto de canto, sus dos secciones a
% distancia R del centro): lo único que cambia es qué lleva el aro, corriente o
% carga. Con corriente las líneas son CERRADAS y atraviesan la espira, de modo
% que en el centro se suman; con carga son radiales y en el centro se cancelan.
% Sentido: sección derecha entrante (x), izquierda saliente (.); para la
% entrante, dB ∝ ds x r con ds = -z y r = -x da +y, y la saliente da también +y.
\begin{center}
\begin{tikzpicture}[scale=1]

  % =============== (a) espira con corriente ===============
  \begin{scope}
    % líneas de campo: cerradas, encerrando cada sección del cable
    \foreach \rx/\ry in {0.45/0.7, 0.82/1.3, 1.15/1.95} {
      \draw[figline, line width=0.7pt] (1.2,0) ellipse ({\rx} and {\ry});
      \draw[figline, line width=0.7pt] (-1.2,0) ellipse ({\rx} and {\ry});
    }
    \draw[figvec, line width=0.7pt] (1.1,1.3) -- (1.32,1.3);
    \draw[figvec, line width=0.7pt] (-1.32,1.95) -- (-1.1,1.95);
    \draw[figvec, line width=0.7pt] (1.32,-1.95) -- (1.1,-1.95);
    \draw[figvec, line width=0.7pt] (-1.1,-1.3) -- (-1.32,-1.3);

    % el aro, visto de canto
    \node[figred, scale=1.15] at (1.2,0) {$\otimes$};
    \node[figred, scale=1.15] at (-1.2,0) {$\odot$};
    \node[figlblaux, figred, anchor=north] at (1.2,-0.2) {$I$};

    % el campo en el centro
    \draw[figvecres, line width=1.4pt] (0,-0.45) -- (0,0.55);
    \node[figlbl, figred, anchor=west] at (0.1,0.2) {$\vec B$};
    \node[figgray, circle, fill=figgray, inner sep=1pt] at (0,-0.45) {};

    \node[figlbl, anchor=north, align=center] at (0,-2.45)
      {(a) espira con corriente:\\[-0.2em]
       $B(0)=\dfrac{\mu_0 I}{2R}\neq 0$};
  \end{scope}

  % =============== (b) el aro cargado, para comparar ===============
  \begin{scope}[xshift=6.4cm]
    % campo radial de cada sección
    \foreach \ang in {30,60,90,120,150,210,240,270,300,330} {
      \draw[figvecaux, figblue] (1.2,0) -- ($(1.2,0)+(\ang:1.15)$);
      \draw[figvecaux, figblue] (-1.2,0) -- ($(-1.2,0)+(\ang:1.15)$);
    }

    % el aro cargado, visto de canto
    \node[figpos, minimum size=7pt] at (1.2,0) {};
    \node[figpos, minimum size=7pt] at (-1.2,0) {};
    \node[figlblaux, figred, anchor=south] at (1.2,1.3) {$+q$};

    % en el centro los aportes se cancelan
    \draw[figvecres, line width=1.2pt] (0,0) -- (-0.75,0);
    \draw[figvecres, line width=1.2pt] (0,0) -- (0.75,0);
    \node[figlbl, anchor=north, align=center] at (0,-1.5)
      {en el centro, aportes opuestos};
    \node[figgray, circle, fill=figgray, inner sep=1pt] at (0,0) {};

    \node[figlbl, anchor=north, align=center] at (0,-2.45)
      {(b) aro cargado:\\[-0.2em] $E(0)=0$};
  \end{scope}

\end{tikzpicture}\\[0.5em]
{\small\itshape La comparación con el aro cargado es el mejor control de que la
cuenta está bien. En el caso eléctrico los aportes en el centro son radiales y
apuntan desde cada trocito hacia el centro: los de secciones opuestas se
cancelan de a pares y queda $E(0)=0$. En el magnético no puede pasar lo mismo
porque $d\vec B$ no es paralelo a $\vec r$ sino perpendicular a $\vec r$ y a
$d\vec s$: los dos aportes se suman y da $B(0)=\mu_0 I/2R$, el máximo de la
curva. La figura muestra además la propiedad estructural del campo magnético:
sus líneas son \textbf{cerradas} ---entran por una cara de la espira y salen por
la otra---, no nacen ni mueren en ninguna parte, porque no hay monopolos
magnéticos. Ese mismo hecho fija el comportamiento lejano: para $Z\gg R$ queda
$B_z\simeq \mu_0 I R^2/2|Z|^3$, un \textbf{dipolo} de momento $\mu = IA = I\pi
R^2$ que decae como $1/r^3$, y no como el $1/r^2$ de una carga puntual. Si se
arreglan dos espiras opuestas para que el dipolo se cancele, lo que sobrevive es
el cuadrupolo, con $1/r^4$.}
\end{center}
